\chapter{Stratégies actions propres à l'Asie}\label{ch:s1:asia-specific-equity-strategies}

Une action de Shanghai atteint sa hausse maximale de dix pour cent à 10:05 et y reste toute la journée~; les acheteurs qui n'ont pas pu négocier sont encore là le lendemain, et elle
ouvre plus haut. Seasholes et Wu ont montré que de tels événements attirent des investisseurs particuliers qui ne détenaient pas le titre, que les cours montent puis reviennent à la moyenne au cours de la
semaine suivante, et que les investisseurs avisés qui achetaient le jour de la limite et revendaient le lendemain gagnaient en moyenne 1.16~\% par jour. Chen, Gao, He, Jiang et Xiong
ont constaté que de gros investisseurs faisaient de même, en poussant les cours jusqu'à la limite pour y parvenir. Sur le marché synthétique à limites de ce chapitre, un acheteur à chaque
clôture en hausse maximale gagnerait 4.5~\% jusqu'à l'ouverture suivante~; un acheteur qui se place dans la file à la limite gagne 2.6~\%, car la file se remplit lorsque le verrouillage est le plus faible~; et
le cours rend 2.8~\% au cours de la semaine suivante. La construction utilise \texttt{firm.limitmkt}.

\section{Limites de variation et T+1}

Le livre~1, chapitre~12 a décrit les règles~: les actions de la cote principale continentale peuvent varier de 10~\% par jour, celles des marchés de croissance de 20~\%, et les actions achetées ne peuvent pas être
revendues le jour même (son encadré daté donne les largeurs en vigueur). Il a aussi montré ce que les limites font aux données~: une limite répartit un grand mouvement sur plusieurs jours, de sorte que
les rendements observés après un jour de limite se prolongent, et un signal qui détecte cette prolongation dans les cours de clôture a trouvé un rendement qui ne peut pas être acheté à
ces cours. C'est sur les transactions qui, elles, sont possibles que porte ce chapitre.

\begin{definition}[Continuation après la limite de variation]\label{def:s1:asia-specific-equity-strategies:continuation}
La \emph{continuation après la limite de variation}\index{continuation après la limite de variation} est la tendance d'une action qui clôture bloquée à sa limite quotidienne à poursuivre son mouvement dans
le même sens à l'ouverture suivante, parce que la part du mouvement retenue par la limite, ainsi que les ordres placés en file à la limite, se reportent au
jour suivant.
\end{definition}

\texttt{firm.limitmkt} produit un marché à limites à partir de rendements équitables (\cref{lst:s1:asia-specific-equity-strategies:limits}). Il prend les dix années
du marché synthétique, sur 1\,000 valeurs, multiplie leurs rendements par 1.5 pour obtenir une volatilité quotidienne de 3.1~\%, et ramène chaque clôture vers la fair value augmentée d'une
prime d'attention, sans jamais dépasser la limite. Une clôture à la limite haute attire le lendemain matin des acheteurs d'attention, qui ajoutent 3~\% à la cible, une prime
qui décroît de 0.6 par jour. La part du mouvement qui dépasse la limite attend l'ouverture suivante. Et un ordre d'achat qui se place dans la file à une limite bloquée est exécuté
avec une probabilité qui baisse à mesure que la valeur de l'action s'éloigne au-delà de la limite~: un verrouillage à peine tenu est exécuté, un verrouillage à fort reliquat ne l'est pas.
Cette dernière règle est la sélection adverse de la file~: les acheteurs les plus susceptibles d'être exécutés sont ceux qui le souhaitaient le moins.

\begin{definition}[Effet d'aimant]\label{def:s1:asia-specific-equity-strategies:magnet}
L'\emph{effet d'aimant}\index{effet d'aimant} est la tendance des cours à accélérer vers une limite de variation de cours quotidienne à mesure qu'ils s'en approchent, plus fortement
vers la limite haute que vers la limite basse, car les intervenants se hâtent de négocier avant que la limite ne les arrête ou ne pousse le cours à se bloquer.
\end{definition}

Cho, Russell, Tiao et Tsay ont documenté l'\omterm{def:s1:asia-specific-equity-strategies:magnet}{effet d'aimant} à Taïwan~; Chen et ses coauteurs l'ont expliqué par de gros investisseurs qui poussent les cours jusqu'à la limite haute
et vendent le lendemain. Le modèle en reproduit le mécanisme~: la moitié des cibles situées à moins de 2~\% sous la limite haute sont poussées jusqu'au blocage, et clôturent au-dessus
de leur fair value. La \cref{fig:s1:asia-specific-equity-strategies:magnet} montre la trace de cet effet~: avec l'aimant, les journées qui auraient clôturé entre 7.5~\% et 9.5~\% de hausse
sont deux fois moins fréquentes, et les clôtures en hausse maximale passent de 0.87~\% des jours-titres à 1.29~\%.

\begin{omfigure}
\begin{tikzpicture}
\begin{axis}[width=10cm,height=5.2cm,ybar=0pt,bar width=6pt,xlabel={\small variation logarithmique de clôture à clôture (\%)},ylabel={\small pour 10\,000 jours-titres},
  tick label style={font=\footnotesize},xmin=4.9,xmax=9.6,ymin=0,ymax=110,grid=major,grid style={gray!25},legend pos=north east,
  legend style={font=\scriptsize},area legend,table/col sep=comma]
\addplot[fill=omExo!50,draw=omExo] table[x=bin_pct,y=none]{figdata/strategies-1/18-asia-specific-equity-strategies/magnet.csv};
\addplot[fill=omDef!70,draw=omDef] table[x=bin_pct,y=magnet]{figdata/strategies-1/18-asia-specific-equity-strategies/magnet.csv};
\legend{sans aimant,avec l'aimant}
\end{axis}
\end{tikzpicture}

\omcaption{Variations quotidiennes juste en deçà d'une limite de 10~\% sur le marché synthétique à limites, par classes d'un demi-point, avec et sans l'aimant (la moitié des cibles situées
à moins de 2~\% de la limite, poussées jusqu'au blocage). Les journées poussées sortent des classes à partir de 7.5~\% et deviennent des clôtures en hausse maximale~: 129 pour 10\,000 jours-titres avec l'aimant, 87
sans. Données~: \texttt{s1\_asia.band}.}\label{fig:s1:asia-specific-equity-strategies:magnet}
\end{omfigure}

Le tableau suivant suit chaque clôture en hausse maximale pendant une semaine (\cref{lst:s1:asia-specific-equity-strategies:ups}). La transaction est celle qu'autorisent les règles~:
acheter à la limite le jour du blocage, vendre à l'ouverture suivante. C'est aussi ce que rend nécessaire la règle T+1, puisque les actions ne peuvent pas être revendues le jour même.

\begin{center}
\small
\begin{tabular}{@{}p{6.0cm}c>{\centering\arraybackslash}p{3.0cm}c@{}}
\toprule
clôtures en hausse maximale à une limite de 10~\% & ensemble & valeur au-delà de la limite & poussées\\
\midrule
part des clôtures & 100~\% & 67~\% & 33~\%\\
rendement jusqu'à l'ouverture suivante, toutes les clôtures & 4.51~\% & 5.90~\% & 1.67~\%\\
probabilité d'exécution dans la file & 57~\% & 36~\% & 100~\%\\
rendement jusqu'à l'ouverture suivante, pondéré par les exécutions & 2.64~\% & 3.96~\% & 1.67~\%\\
de l'ouverture suivante à la clôture cinq jours plus tard & $-2.81~\%$ & $-2.77~\%$ & $-2.88~\%$\\
\bottomrule
\end{tabular}
\end{center}

Deux effets font le gain de nuit. Le reliquat est la part du mouvement retenue par la limite~; sans attention ni aimant, il rapporte 3.42~\% jusqu'à
l'ouverture suivante (1.21~\% après exécutions). Les acheteurs d'attention ajoutent l'essentiel du reste, et le reprennent au cours de la semaine. Les verrouillages poussés rapportent peu pendant la nuit, car
ils ont clôturé au-dessus de leur valeur, mais ils attirent tout de même les acheteurs d'attention. Les exécutions réduisent le gain de 4.51~\% à 2.64~\%~: l'acheteur qui n'a pas pu être exécuté dans
les verrouillages les plus forts est la règle, non l'exception. Pour un aller-retour de 15~points de base (le droit de timbre à la vente de l'encadré daté, plus des commissions
et une fourchette supposées), la transaction exécutée dégage encore 2.49~\% en moyenne, soit environ le double du 1.16~\% des investisseurs avisés chez Seasholes et Wu, car le reliquat et
l'attention du modèle sont propres et chaque position dans la file est équivalente.

\begin{omfigure}
\begin{tikzpicture}
\begin{axis}[width=10cm,height=5.4cm,ybar,bar width=10pt,xlabel={\small limite de variation de cours quotidienne (\%)},ylabel={\small rendement moyen (\%)},
  tick label style={font=\footnotesize},symbolic x coords={5,10,20},xtick=data,ymin=-6,ymax=8,enlarge x limits=0.25,grid=major,
  grid style={gray!25},legend columns=1,legend style={font=\scriptsize,at={(0.5,-0.3)},anchor=north,draw=none,fill=none,
  /tikz/every even column/.append style={column sep=8pt}},area legend,table/col sep=comma]
\addplot[fill=omThm!60,draw=omThm] table[x=limit_pct,y=open_pct]{figdata/strategies-1/18-asia-specific-equity-strategies/widths.csv};
\addplot[fill=omDef!60,draw=omDef] table[x=limit_pct,y=filled_pct]{figdata/strategies-1/18-asia-specific-equity-strategies/widths.csv};
\addplot[fill=omExo!50,draw=omExo] table[x=limit_pct,y=week_pct]{figdata/strategies-1/18-asia-specific-equity-strategies/widths.csv};
\legend{de la clôture à l'ouverture suivante,pondéré par les exécutions dans la file,de l'ouverture suivante à cinq jours plus tard}
\end{axis}
\end{tikzpicture}

\omcaption{Clôtures en hausse maximale sur le marché synthétique à limites, selon la largeur de la limite quotidienne de variation~: le rendement moyen de la clôture à l'ouverture suivante pour toutes les clôtures
et pondéré par la probabilité qu'un ordre d'achat placé en file soit exécuté, et le rendement moyen sur la semaine suivante. Données~:
\texttt{s1\_asia.limit\_ups}.}\label{fig:s1:asia-specific-equity-strategies:widths}
\end{omfigure}

La largeur décide de la fréquence et de l'intensité (\cref{fig:s1:asia-specific-equity-strategies:widths}). Une limite de 5~\% se bloque sur 11.24~\% des jours-titres (les
actions synthétiques sont volatiles), une limite de 10~\% se bloque sur 1.29~\%, et une limite de 20~\% sur 0.10~\%. Plus la limite est large, plus le mouvement qui précède chaque blocage est extrême~: le
gain à l'ouverture suivante monte de 2.98~\% à 4.51~\% puis à 6.58~\%, le gain après exécutions de 2.12~\% à 2.64~\% puis à 3.20~\%, et le recul de la semaine de 1.76~\% à 2.81~\% puis
à 5.27~\%. La colonne à 20~\% repose sur 2\,325 événements, et son chiffre de la semaine suivante est bruité~; les moyennes à l'ouverture suivante ont des erreurs types de 0.02 point à 10~\% et
0.1 point à 20~\%.

\begin{dated}{2026-09}{Coûts en Chine continentale et données Stock Connect}\label{dat:s1:asia-specific-equity-strategies:connect}
À partir du 28 août 2023, le droit de timbre sur les transactions d'actions de Shanghai et de Shenzhen a été réduit de moitié, de 0.1~\% à 0.05~\% du montant de la transaction, payé par le
vendeur. Le 12 avril 2024, les bourses de Hong Kong, de Shanghai et de Shenzhen ont annoncé des modifications des données Northbound de Stock Connect, en deux phases, environ un
et quatre mois plus tard~: les montants négociés en temps réel à l'achat, à la vente et au total ne seraient plus disponibles (les totaux quotidiens et mensuels ainsi que les dix actions les plus actives
le seraient)~; le quota quotidien restant ne serait affiché que lorsqu'il est inférieur à 30~\%~; et les participations Northbound seraient publiées chaque trimestre, le cinquième
jour de négociation Northbound suivant la fin du trimestre.
\end{dated}

\section{Marchés dominés par les particuliers}

La majeure partie des négociations sur les actions continentales est le fait de particuliers, et ceux-ci ne forment pas un seul groupe. Jones, Shi, Zhang et Zhang répartissent les comptes de particuliers chinois selon
taille~: les plus petits ne parviennent pas à prévoir les rendements, courent après les mouvements du jour, n'intègrent pas les nouvelles publiques et manifestent des préférences pour le jeu~; les plus gros prévoient
correctement les rendements, et opèrent à contre-courant des mouvements du jour, bien que leurs coûts absorbent ce qu'ils gagnent. Les achats d'attention après les jours de limite haussière sont le
comportement des petits comptes, et les grands investisseurs de Chen et coauteurs se trouvent de l'autre côté.

La semaine suivant une clôture à la limite haussière est le terrain de la stratégie des particuliers. Sur le marché synthétique, la prime d'attention s'inverse de 2.81~\% depuis l'ouverture suivante
cinq jours plus tard. En profiter suppose une vente à découvert, et l'emprunt de titres peut être rare sur ces marchés ou, comme en Corée avec ses dix-sept mois sans
vente à découvert (livre~1, chapitre~12), purement et simplement interdit. En pratique, on se trouve de l'autre côté~: un détenteur vend sur la vague d'attention à l'ouverture, et un portefeuille qui aurait
acheté dans le prolongement reste à l'écart.

\section{Flux Connect}

\begin{definition}[Flux Northbound]\label{def:s1:asia-specific-equity-strategies:northbound}
Le \emph{flux Northbound}\index{flux Northbound} est l'achat net d'actions continentales par des investisseurs qui négocient via Stock Connect depuis Hong Kong, par
action ou en agrégé~; le \emph{flux Southbound} est l'inverse, les investisseurs continentaux achetant des actions de Hong Kong.
\end{definition}

Les investisseurs Northbound sont surtout des institutions, et Bian, Chan, Han et Shi ont constaté que leurs achats nets prévoient les rendements des actions de Shanghai
connectées. Une stratégie qui les suit a besoin de leurs flux action par action, et vite~: l'information contenue dans un flux s'épuise rapidement. L'encadré daté indique que les
participations par action, publiées chaque jour à l'origine, sont désormais trimestrielles. Le chapitre introduit des flux qui connaissent un peu les cinq jours suivants de chaque action (une corrélation de
0.02 avec le rendement) et construit sur eux un portefeuille long--short hebdomadaire, à 15 points de base par aller-retour.

\begin{center}
\small
\begin{tabular}{@{}lc>{\centering\arraybackslash}p{4.2cm}@{}}
\toprule
participations Northbound publiées & quotidiennes & trimestrielles, avec cinq jours de retard\\
\midrule
IC avec le rendement des cinq jours suivants & 0.021 & 0.003\\
ratio de Sharpe après coûts, portefeuille hebdomadaire & 2.04 & 0.26\\
\bottomrule
\end{tabular}
\end{center}

Avec des données quotidiennes, l'IC de 0.021 du flux, sur un millier de valeurs par semaine, donne un ratio de Sharpe de 2.0 après coûts~; avec des données trimestrielles publiées avec cinq jours de retard,
les mêmes flux ne valent presque rien. Le faible IC qui subsiste (0.003, avec une erreur type de 0.001) provient des dérives
persistantes du marché synthétique, qui rendent les achats informés du trimestre précédent faiblement informatifs pour la semaine suivante. La leçon est générale~: un signal fait des
opérations publiées d'autrui vaut ce que le calendrier de publication lui laisse.

\section{Fiches de stratégie}

\begin{strategyfile}{Continuation après la limite de variation}\label{strat:s1:asia-specific-equity-strategies:limit}
\sfield{Qui paie, et pourquoi} Acheteurs bloqués par la limite, et acheteurs d'attention qui arrivent le lendemain matin.
\sfield{Instruments et plateformes} Actions A continentales à la limite quotidienne~; autres marchés avec des limites quotidiennes.
\sfield{Signal} Une clôture à la limite haussière, de préférence avec un important arriéré d'ordres d'achat non exécutés et adossée à une nouvelle.
\sfield{Dimensionnement et exécution} Un ordre d'achat placé tôt dans la file à la limite, revendu à l'ouverture suivante (la règle T+1 n'autorise pas plus tôt).
\sfield{Coûts} Droit de timbre sur la vente, commissions~; la sélection adverse dans la file, qui constitue le principal coût.
\sfield{Comment elle s'éteint} Concurrence pour la priorité dans la file~; l'attention des régulateurs aux gros acheteurs à la limite.
\sfield{Horizon, capacité, infrastructure} Une nuit~; capacité limitée par ce que la file permet d'exécuter~; transmission des ordres assez rapide pour arriver tôt dans la file.
\sfield{Backtest en toute honnêteté} Exécutions issues de la file, et non des cours de clôture~; une clôture à la limite n'est pas un prix auquel quiconque aurait pu acheter.
\sfield{Sources} Seasholes et Wu (2007)~: 1.16~\% par jour pour les investisseurs avisés~; Chen, Gao, He, Jiang et Xiong (2019).
\end{strategyfile}

\begin{strategyfile}{Retournement de l'attention des particuliers}\label{strat:s1:asia-specific-equity-strategies:retail}
\sfield{Qui paie, et pourquoi} Petits comptes de particuliers qui achètent des actions captant l'attention après les jours de limite haussière.
\sfield{Instruments et plateformes} Actions après les jours de limite haussière~; pour une vente à découvert, celles qui peuvent être empruntées.
\sfield{Signal} Une clôture à la limite haussière la veille, l'écart à l'ouverture, la part des particuliers dans la négociation.
\sfield{Dimensionnement et exécution} Vendre les titres détenus ou vendre à découvert à l'ouverture~; conserver la position une semaine.
\sfield{Coûts} Coût de l'emprunt, là où il est disponible~; droit de timbre.
\sfield{Comment elle s'éteint} Restrictions sur les ventes à découvert~; moins de négociation par les particuliers.
\sfield{Horizon, capacité, infrastructure} Une semaine~; capacité limitée par l'emprunt.
\sfield{Backtest en toute honnêteté} Seulement les actions empruntables le jour même~; le cours d'ouverture, et non le cours de clôture de la veille.
\sfield{Sources} Seasholes et Wu (2007)~: retour à la moyenne sur la semaine suivante~; Jones, Shi, Zhang et Zhang (2025).
\end{strategyfile}

\begin{strategyfile}{Signal de flux Northbound}\label{strat:s1:asia-specific-equity-strategies:northbound}
\sfield{Qui paie, et pourquoi} Investisseurs plus lents que les institutions qui négocient via Stock Connect, et qui apprennent leur information plus tard.
\sfield{Instruments et plateformes} Actions continentales éligibles à Stock Connect.
\sfield{Signal} Variations des participations Northbound par action, ou \omterm{def:s1:asia-specific-equity-strategies:northbound}{flux Northbound} agrégé.
\sfield{Dimensionnement et exécution} Un portefeuille long--short sur les variations, rééquilibré aussi souvent que les données le permettent.
\sfield{Coûts} Droit de timbre et commissions à chaque rééquilibrage.
\sfield{Comment elle s'éteint} Les données~: les participations par action sont désormais publiées chaque trimestre.
\sfield{Horizon, capacité, infrastructure} De quelques jours à quelques semaines avec des données quotidiennes~; le calendrier des données le fixe.
\sfield{Backtest en toute honnêteté} Participations telles que publiées le jour de leur publication, avec le calendrier en vigueur à l'époque.
\sfield{Sources} Bian, Chan, Han et Shi (2023)~; l'encadré daté.
\end{strategyfile}

\section{Tutoriel~: limite haussière verrouillée}\label{tut:s1:asia-specific-equity-strategies:tut}

\begin{tutorial}
\textbf{Objectif.} Reproduire le marché synthétique sous limites quotidiennes de variation de cours, suivre chaque clôture à la limite haussière et mesurer ce qu'auraient obtenu un acheteur à la limite, un acheteur placé
dans la file et un vendeur une semaine plus tard. \textbf{État final~:} les tableaux et les deux figures.
\begin{tutsteps}
\item \textbf{La couche des limites}~: ouvertures, clôtures, verrouillages, \omterm{def:s1:asia-specific-equity-strategies:magnet}{effet d'aimant} et exécutions dans la file, jour après jour.
\omcode{code/firm/limitmkt/firm_limitmkt.py}{43}{73}{Reproduction d'un marché à limites à partir des rendements de fair value.}\label{lst:s1:asia-specific-equity-strategies:limits}
\item \textbf{Les événements}~: chaque clôture à la limite haussière suivie d'une semaine de cotation, réparties entre verrouillages avec arriéré et verrouillages poussés.
\omcode{code/strategies-1/18-asia-specific-equity-strategies/python/s1_asia.py}{47}{64}{Clôtures à la limite haussière et ce qui les suit.}\label{lst:s1:asia-specific-equity-strategies:ups}
\item \textbf{Lancer} \texttt{limit\_ups(limit)} pour les trois largeurs et les deux témoins (\texttt{attention=0}, \texttt{push=0}), \texttt{band(push)},
\texttt{flows(1, 0)} et \texttt{flows(63, 5)}, ainsi que \texttt{fig\_asia.py}.
\end{tutsteps}
\textbf{Pour aller plus loin.} Introduire des positions dans la file~: les premiers acheteurs sont exécutés en premier~; faire dépendre la prime d'attention du nombre d'actions qui se verrouillent le même jour
(Seasholes et Wu ont constaté moins d'acheteurs lorsque beaucoup le font)~; ajouter des verrouillages à la limite basse et les vendeurs qu'ils piègent.
\end{tutorial}

\section{Construction~: un marché à limites}\label{bld:s1:asia-specific-equity-strategies:limitmkt}

\begin{build}
\bfield{Rôle} Prix d'un marché à limites quotidiennes à partir de rendements fair value, avec achats d'attention, \omterm{def:s1:asia-specific-equity-strategies:magnet}{effet d'aimant} et exécutions en file~; un générateur de \omterm{def:s1:asia-specific-equity-strategies:northbound}{flux Northbound} et
son calendrier de publication.
\bfield{Interface} \texttt{LimitConfig(\dots)}, \texttt{apply\_limits(R, cfg, rng)}, \texttt{northbound(alpha, skill, rng)},
\texttt{holdings\_seen(flow, every, lag)}.
\bfield{Règles} Aucune clôture au-delà de la limite fixée depuis la clôture précédente~; l'ouverture reporte l'arriéré de la veille et la prime du jour, dans cette limite~; une cotation
démarre à sa fair value.
\bfield{Tests d'acceptation} \texttt{code/firm/limitmkt/tests/}~: un mouvement de 15~\% se verrouille à 10~\% avec la probabilité d'exécution calculée à la main et ouvre à la fair value
plus la prime~; l'\omterm{def:s1:asia-specific-equity-strategies:magnet}{effet d'aimant} tire vers une cible de 9~\%~; les participations publiées tous les quatre jours avec deux jours de retard~; la corrélation des flux.
\bfield{Pour aller plus loin} Priorité dans la file~; verrouillages à la limite basse~; une contrainte T+1 sur les allers-retours intrajournaliers dans le backtester.
\end{build}

\begin{omsources}
\item T.~Chen, Z.~Gao, J.~He, W.~Jiang et W.~Xiong, «~Daily price limits and destructive market behavior~», \emph{Journal of Econometrics} 208(1),
2019.
\item D.~D.~Cho, J.~Russell, G.~C.~Tiao et R.~Tsay, «~The magnet effect of price limits: evidence from high-frequency data on Taiwan Stock Exchange~»,
\emph{Journal of Empirical Finance} 10(1--2), 2003.
\item M.~S.~Seasholes et G.~Wu, «~Predictable behavior, profits, and attention~», \emph{Journal of Empirical Finance} 14(5), 2007.
\item C.~M.~Jones, D.~Shi, X.~Zhang et X.~Zhang, «~Retail trading and return predictability in China~», \emph{Journal of Financial and Quantitative
Analysis} 60(1), 2025 (en ligne 2024).
\item J.~Bian, K.~Chan, B.~Han et D.~Shi, «~Cross-border equity flows and information transmission: evidence from Chinese stock markets~»,
\emph{Journal of International Financial Markets, Institutions and Money} 84, 2023.
\end{omsources}

\section{Exercices}

\begin{exercise}[$\star$]\label{exo:s1:asia-specific-equity-strategies:1}
L'actualité fait monter la valeur d'une action de 50~\%. Combien de jours faut-il pour l'atteindre avec une limite de 10~\%, et avec une limite de 20~\%~?
\end{exercise}

\begin{exercise}[$\star$]\label{exo:s1:asia-specific-equity-strategies:2}
La valeur d'une action progresse de 15~\% en une journée avec une limite de 10~\%. Avec une échelle d'exécution de la couche de 2~\%, quelle est la probabilité qu'un ordre d'achat mis en file soit exécuté~?
\end{exercise}

\begin{exercise}[$\star$]\label{exo:s1:asia-specific-equity-strategies:3}
La prime d'attention vaut 3~\% à l'ouverture qui suit une clôture sur la limite haute et décroît de 0.6 par jour. Quelle part en subsiste quatre jours plus tard, et quelle part
s'est résorbée~?
\end{exercise}

\begin{exercise}[$\star\star$]\label{exo:s1:asia-specific-equity-strategies:4}
Pourquoi les verrouillages poussés ont-ils une probabilité d'exécution égale à un et un faible gain de nuit~?
\end{exercise}

\begin{exercise}[$\star\star$]\label{exo:s1:asia-specific-equity-strategies:5}
Pourquoi le gain de nuit croît-il avec la largeur de la limite~?
\end{exercise}

\begin{exercise}[$\star\star$]\label{exo:s1:asia-specific-equity-strategies:6}
Avec une limite de 10~\%, 1.29~\% des jours-action sont des clôtures à la limite haute. Combien un marché de 1\,000 actions en produit-il par jour, et combien rapporte chaque transaction exécutée
sur chacune, après un aller-retour de 15 points de base~?
\end{exercise}

\begin{exercise}[$\star\star\star$]\label{exo:s1:asia-specific-equity-strategies:7}
\emph{Programmation.} Exécuter \texttt{limit\_ups(0.10, attention=0.0, push=0.0)}. Que reste-t-il du gain de nuit, exécuté et non exécuté, et du
retournement de la semaine~? Qu'est-ce que cela indique sur l'origine du rendement de la transaction~?
\end{exercise}

\begin{exercise}[$\star\star\star$]\label{exo:s1:asia-specific-equity-strategies:8}
\emph{Trouver la faille.} «~Les actions qui clôturent à la limite haute gagnent 4.5~\% à l'ouverture suivante en moyenne~; acheter chaque clôture à la limite haute revient à obtenir un ratio de Sharpe de dix.~»
\end{exercise}

\section{Problème~: Verrouillage à la limite haute}

\begin{problem}[{Problème du week-end --- une journée à la limite}]\label{pb:s1:asia-specific-equity-strategies:1}
Le marché synthétique à limites du chapitre et les travaux publiés.

\medskip\noindent\textbf{Partie I --- Les règles.}
\begin{enumerate}
\item Définir la \omterm{def:s1:asia-specific-equity-strategies:continuation}{continuation après la limite de variation} et l'\omterm{def:s1:asia-specific-equity-strategies:magnet}{effet d'aimant}.
\item Quel effet les limites quotidiennes ont-elles sur les rendements observés (livre~1, chapitre~12)~?
\item Qu'interdit la règle T+1, et quelle transaction laisse-t-elle encore~?
\item Décrire la couche de limites du chapitre.
\end{enumerate}

\medskip\noindent\textbf{Partie II --- Les travaux publiés.}
\begin{enumerate}[resume]
\item Qu'ont établi Seasholes et Wu après les événements de limite haute~?
\item Qu'ont établi Chen et ses coauteurs sur les gros investisseurs~?
\item Qu'ont établi Jones et ses coauteurs sur les comptes de particuliers de tailles différentes~?
\item Qu'ont établi Bian et ses coauteurs sur les \omterm{def:s1:asia-specific-equity-strategies:northbound}{flux Northbound}~?
\end{enumerate}

\medskip\noindent\textbf{Partie III --- Les événements.}
\begin{enumerate}[resume]
\item À quelle fréquence les clôtures à la limite haute surviennent-elles avec des limites de 5~\%, 10~\% et 20~\%~?
\item Donner le gain de nuit pour chaque clôture, puis pondéré par les exécutions, avec une limite de 10~\%.
\item Pourquoi les exécutions réduisent-elles le gain~?
\item Comment l'\omterm{def:s1:asia-specific-equity-strategies:magnet}{effet d'aimant} modifie-t-il la distribution des variations quotidiennes~?
\end{enumerate}

\medskip\noindent\textbf{Partie IV --- Le verdict.}
\begin{enumerate}[resume]
\item Énoncer le \emph{résultat nommé}~: le rendement du lendemain après un verrouillage à la limite haute et la manière dont il varie avec la largeur de la limite.
\item Décomposer le gain de nuit entre arriéré et attention.
\item Pourquoi le retournement lié à l'attention des particuliers est-il difficile à exploiter en Chine continentale~?
\item Quel effet le changement de données de 2024 a-t-il eu sur le signal des \omterm{def:s1:asia-specific-equity-strategies:northbound}{flux Northbound}~?
\item Comment backtester honnêtement une stratégie de limite haute~?
\item Quelle fiche de stratégie dépend le plus de la réglementation~?
\item Quel lien ce chapitre entretient-il avec le chapitre~9 sur les flux~?
\item En une phrase~: que vend une limite de variation de cours, et à qui~?
\end{enumerate}
\end{problem}

\section{Questions d'entretien}

\begin{interviewq}[$\star$ \iqroles{chercheur}]\label{iq:s1:asia-specific-equity-strategies:1}
Pourquoi un signal de momentum calibré sur les cours de clôture d'un marché à limites quotidiennes est-il suspect~?
\end{interviewq}

\begin{interviewq}[$\star\star$ \iqroles{trader}]\label{iq:s1:asia-specific-equity-strategies:2}
Une action que vous voulez acheter est verrouillée à la limite haute avec une file importante. Que faites-vous~?
\end{interviewq}

\begin{interviewq}[$\star\star$ \iqroles{chercheur}]\label{iq:s1:asia-specific-equity-strategies:3}
Comment détecteriez-vous l'\omterm{def:s1:asia-specific-equity-strategies:magnet}{effet d'aimant} dans les données intrajournalières~?
\end{interviewq}

\begin{interviewq}[$\star\star$ \iqroles{développeur}]\label{iq:s1:asia-specific-equity-strategies:4}
Qu'est-ce qui change dans un backtester pour un marché à limites quotidiennes et la règle T+1~?
\end{interviewq}

\begin{interviewq}[$\star\star$ \iqroles{risque}]\label{iq:s1:asia-specific-equity-strategies:5}
Quels sont les risques liés à la détention d'un portefeuille d'actions de Chine continentale susceptibles d'être verrouillées à la limite basse~?
\end{interviewq}

\begin{interviewq}[$\star\star\star$ \iqroles{chercheur}]\label{iq:s1:asia-specific-equity-strategies:6}
La valeur d'une action bondit de $x$ au-dessus d'une limite $L$ (en termes logarithmiques). Les ordres d'achat se mettent en file à la limite et sont exécutés avec la probabilité $e^{-(x - L)/s}$ pour $x > L$~;
l'ouverture suivante reflète la valeur complète. Déterminer le gain de nuit espéré d'un ordre exécuté lorsque $x - L$ suit une loi exponentielle de moyenne $m$, et
le comparer à la moyenne inconditionnelle.
\end{interviewq}
